% Experimental setup section — drafted 2026-07-11 for the ICASP 2026 paper.
% Numbers match the current repo state (icasp_paper.ipynb + scripts/).
% Figure references left as \ref{...} placeholders — wire them to the Overleaf labels.

\section{Experimental Setup}
\label{sec:setup}

\subsection{Protocol}
All experiments compare three aggregation rules under an identical
availability-driven participation process. There are $N$ clients with local
objectives $F_i$, an \emph{importance} distribution $p$ defining the global
objective $F(w)=\sum_{i=1}^{N} p_i F_i(w)$, and an \emph{availability}
distribution $r$ governing participation. In every round the server draws a
subset $S$ of $K$ clients without replacement with probabilities proportional
to $r$ (implemented via the Gumbel top-$K$ trick). Each selected client runs
$E=5$ epochs of full-batch gradient descent from the current global model,
and the three rules aggregate the resulting local models $\theta_i$ in
parallel from the same subset draw, so their loss curves are directly
comparable:
\begin{itemize}
  \item \textbf{FedAVOT}: $w \leftarrow \sum_{i \in S} T_{ij}\,\theta_i$,
        where $j$ indexes the sampled subset and $T$ is the masked
        optimal-transport plan described below;
  \item \textbf{FedAvg($K$)} (partial participation):
        $w \leftarrow \sum_{i \in S} \tfrac{N}{K}\, p_i\, \theta_i$;
  \item \textbf{FedAvg(full)} (full-participation reference):
        $w \leftarrow \sum_{i=1}^{N} p_i\, \theta_i$ over all clients.
\end{itemize}

\textbf{Transport plan.} The columns of the transport problem are all
$\binom{N}{K}$ subsets. The column marginal $q_j=\Pr[S_j \text{ drawn}]$ is
estimated by Monte Carlo ($10^6$ draws for the real-data experiments,
$4\times10^5$ for the synthetic ones). A joint coupling $Y$ with row marginal
$p$ and column marginal $q$, masked to subset membership, is fitted by
Sinkhorn-style iterative proportional fitting (IPFP; tolerance $10^{-12}$,
at most $10^3$ iterations for real data; $10^{-10}$ and $3\times10^3$ for
synthetic). Each column of $Y$ is then normalized to a convex weight vector,
$T_{\cdot j} = Y_{\cdot j}/q_j$, giving the FedAVOT aggregation weights.

\textbf{Feasibility.} Because $T$ can place weight on a client only through
subsets containing it, the marginal constraint $\sum_j q_j T_{ij} = p_i$ is
achievable only if $p_i \le \pi_i$ for every client, where
$\pi_i=\Pr[i \in S]$ is the inclusion probability (estimated with
$5\times10^5$ Monte Carlo draws). We quantify infeasibility as the importance
mass held by violating clients, $\sum_{i:\,p_i>\pi_i} p_i$.

\textbf{Metrics.} We report the $p$-weighted global mean-squared error
$\sum_i p_i \,\mathrm{MSE}_i(w)$ per round; final values are averaged over
the last $500$ rounds and reported as mean $\pm$ standard deviation across
seeds.

\subsection{Real data: IMDb-Wiki age regression}
Clients are formed by grouping IMDb-Wiki images by celebrity identity; we
take the first $N=100$ identities with at least $30$ training images and
give each client $30$ images. Features are $128$-dimensional face embeddings
extracted with a pretrained ResNet; the task is linear regression of
(mean-centered) age on the embedding. Hyperparameters: $K=3$, $4000$ rounds
($5000$ for the feasible-availability variant), learning rate $10^{-2}$,
$E=5$ local epochs, $5$ seeds. Two availability regimes are evaluated
against the same cubic importance $p_i \propto (N{+}1{-}i)^3$:
\begin{itemize}
  \item \emph{Mirrored (infeasible)}: $r_i \propto i^3$, the exact mirror
        image of $p$, so the most important clients are almost never
        available ($41/100$ clients violate $p_i \le \pi_i$, holding
        $\approx 88\%$ of the importance mass);
  \item \emph{Aligned (feasible)}: $r_i \propto N{+}1{-}i$, a linear skew in
        the same direction as $p$. The transport is feasible for all clients
        ($\max_i p_i/\pi_i = 0.66$) while the cubic-vs-linear shape mismatch
        still requires nontrivial reweighting.
\end{itemize}

\subsection{Controlled synthetic experiment}
To isolate the effect of feasibility we use a linear-regression setting in
which the skew of both distributions is a single knob $\alpha$:
$p_i \propto (N{+}1{-}i)^{\alpha}$ and $r_i \propto i^{\alpha}$, with
$N=50$, $K=3$. Client $i$ holds $30$ samples with standard-normal features
in $d=40$ dimensions and labels generated by a client-specific linear model
$w_i = \bar w + h\,\delta_i u + 0.3\,\varepsilon_i$ with label noise
$\sigma=0.1$, where $\delta_i = 2i/(N{-}1)-1$ orders the drift along the
availability axis and $h=3$ sets its strength — that is, the distribution
shift is availability-correlated, so which clients participate determines
which optimum is reached. Hyperparameters: learning rate $10^{-3}$, $E=5$,
$1500$ rounds ($5000$ for the extended feasible run), $3$ seeds. We sweep
$\alpha \in \{0.5, 1, 1.5, 2, 3\}$, spanning $\approx 14\%$ to
$\approx 87\%$ infeasible importance mass; the feasibility-mechanism
diagnostic contrasts $\alpha=0.2$ (fully feasible, $0\%$) with $\alpha=3.0$.
